% Research fragment for the parent article.  Not a standalone document.
% Authored against ProveIt f7e35be357654198c1546fd40f363694c5bcc2cc.
% Source conjecture: incoming Canonical Digamma and Fifth Harmonic Bridge,
% sections/09_questions.tex, Conjecture conj:firstq.
% Source exact tail formulas: sections/07_tail_rank.tex, tr:Aentries,
% tr:forward, tr:backward, tr:operator.

\section{The optimal moment prefix for adjacent even tail orders}
\label{adj:section}

The incoming \emph{Canonical Digamma Polynomials and the Fifth Harmonic
Bridge} proves that finitely many initial centered moments generate the
full formal coordinate space of two even harmonic tails.  For equal
orders it proves the optimal consecutive prefix.  Its
Conjecture~\texttt{conj:firstq} asks for that optimal prefix for every
unequal pair.  We prove the entire first off-diagonal family:
\[
 p=2b-2,\qquad q=2b,\qquad b\ge2.
\]
The proof also gives an exact determinant with a positive rational
recurrence, and hence a family of ordinary convergent identities for
products of higher polygammas.  The result concerns rational coefficient
rank.  It does not assert arithmetic independence of the constants that
those coefficients multiply.

\subsection{The inherited exact coordinate system}

Write
\[
 A_p(n)=\zeta(p)-H_n^{(p)}=\zeta(p,n+1),\qquad
 x_n=n+\tfrac12,
\]
and let $D_r$ be the continued value at $s=-2r$ of
\[
 D_{p,q}(s)=\sum_{n=0}^{\infty}A_p(n)A_q(n)x_n^{-s}.
\]
For even $p,q$, these particular continued values are regular.  The
source tail theorem supplies numbers $C_j$ satisfying the inverse
Bernoulli transforms
\begin{align}
 D_r&=\sum_{j=0}^{r}
 \frac{\binom{2r}{2j}}{2j+1}B_{2r-2j}(1/2)C_j,
 \label{adj:forward}\\
 C_j&=\sum_{r=0}^{j}\binom{2j+1}{2r}2^{-2(j-r)}D_r.
 \label{adj:backward}
\end{align}
For $q=2b$ and $j\ge b$, they have the finite coordinate expansion
\begin{equation}\label{adj:block}
 C_j=\sum_{k=0}^{b-1}a_{jk}\zeta(2k),\qquad \zeta(0)=-\tfrac12,
\end{equation}
where, putting $d=2j+1$ and $e=d-p-q+2k+1$,
\begin{equation}\label{adj:a}
 a_{jk}=B_e\left\{
 \frac{\binom{d-q+1}{p-2k}}{d-q+1}
 +\frac{\binom{d-p+1}{q-2k}}{d-p+1}
 \right\}\quad(e\ge0),
\end{equation}
and $a_{jk}=0$ if $e<0$.  Each binomial summand in
\eqref{adj:a} is omitted unless its lower index is at least $2$ and
no greater than its upper index.  In particular, a would-be lower
index zero is \emph{not} retained.  These formulas are prior results
of the inspected reports; the new argument is the nonvanishing theorem
below.

For reference, the finite mechanism behind the inherited formulas is
summation by parts with
\[
 \sum_{n=0}^{N-1}x_n^{2r}
 =\sum_{j=0}^{r}\frac{\binom{2r}{2j}}{2j+1}
 B_{2r-2j}(1/2)N^{2j+1},
 \qquad A_p(k-1)=A_p(k)+k^{-p}.
\]
The resulting ordered finite sums reduce by Faulhaber polynomials.
Their constant coefficients give \eqref{adj:block}--\eqref{adj:a};
all subtracted nonnegative even powers have continued sum zero because
$\zeta(-2m,1/2)=0$.  Thus the conventions in
\eqref{adj:forward} are tied to the original centered cutoff, rather
than to an unspecified regularization.

\subsection{A universal positive recurrence}

Let
\[
 h(t)=\coth(t/2)=\sum_{k=0}^{\infty}h_k t^{2k-1},\qquad
 h_k=\frac{2B_{2k}}{(2k)!},\qquad h_0=2.
\]
The ordinary Bernoulli sign formula gives
\begin{equation}\label{adj:hsign}
 h_k=(-1)^{k-1}\beta_k,\qquad
 \beta_k=\frac{2|B_{2k}|}{(2k)!}
 =\frac{4\zeta(2k)}{(2\pi)^{2k}}>0\quad(k\ge1).
\end{equation}
This is the $x=0$ case of the classical Bernoulli Fourier formula
(DLMF~24.8.1).  Define rational numbers $c_n$ by
\begin{align}
 c_1&=1,\label{adj:c1}\\
 c_n&=\frac{1}{4(2n-1)^2}
 \sum_{r=1}^{n-1}\beta_{n-r}
 \left\{2r(2r-1)+(2n-1)(2n-2)-2\mathbf1_{r=1}\right\}c_r
 \quad(n\ge2).
 \label{adj:positive-recurrence}
\end{align}
Every coefficient in this recurrence is strictly positive.  Indeed,
for $r=1$ the expression in braces is $(2n-1)(2n-2)>0$, and for
$r\ge2$ both its displayed summands are positive.  Therefore
\begin{equation}\label{adj:cpositive}
 c_n\in\mathbb Q_{>0}\quad(n\ge1).
\end{equation}
The first values are
\[
 c_1=1,\quad c_2=\frac1{36},\quad c_3=\frac{11}{5400},\quad
 c_4=\frac{113}{705600},\quad
 c_5=\frac{2509}{190512000},\quad
 c_6=\frac{309713}{276623424000}.
\]

\begin{theorem}[Optimal prefix and exact determinant for adjacent orders]
\label{adj:main}
Let $b\ge2$, $p=2b-2$, and $q=2b$.  Define the square late-block
matrix
\[
 M_b=(a_{b+r,k})_{0\le r,k<b}
\]
by \eqref{adj:a}.  Then
\begin{equation}\label{adj:determinant}
 \det M_b=
 (-1)^{b(b-1)/2}\frac{((2b+1)!!)^2}{b!}\,c_{b+1}\ne0.
\end{equation}
Consequently $D_0,\ldots,D_{q-1}$ have full formal rank $q$ in
the coordinate system
\[
 C_0,\ldots,C_{b-1},\quad\zeta(0),\zeta(2),\ldots,\zeta(2b-2).
\]
Every later $D_r$ is an explicit rational combination of these first
$q$ moments.  In this coordinate ordering, the determinant of the
complete first-$q$ moment matrix is
\begin{equation}\label{adj:full-determinant}
 (-1)^{b(b-1)/2}
 \frac{((2b+1)!!)^2c_{b+1}}{b!(4b-1)!!}.
\end{equation}
\end{theorem}

\begin{proof}
For an even polynomial
\[
 Q(t)=\sum_{n=1}^{b}q_n t^{2n},
\]
introduce the operator
\begin{equation}\label{adj:operator}
 \mathcal L_Q(t)=\bigl(Q''(t)-Q''(0)\bigr)h(t)+(Q(t)h(t))''.
\end{equation}
It is an odd formal power series, and is holomorphic at zero when
regarded analytically near zero.  The subtraction of $Q''(0)$ is
essential: it is precisely the omitted lower-index-zero term in
\eqref{adj:a}.

Let $K_b$ be the matrix mapping $(q_1,\ldots,q_b)$ to the
coefficients of $t,t^3,\ldots,t^{2b-1}$ in $\mathcal L_Q$.
Direct coefficient extraction from \eqref{adj:a} gives
\begin{equation}\label{adj:matrix-conjugacy}
 K_b=\operatorname{diag}_{0\le r<b}\left(\frac2{(2r+1)!}\right)
 M_b^{\mathrm{rev}}
 \operatorname{diag}_{1\le n\le b}\bigl((2n)!\bigr),
\end{equation}
where $M_b^{\mathrm{rev}}$ reverses the column order.  This can
also be read off from the inherited identity
$2f^{(2b)}=\mathcal L_Q$ for the column generators, but no
analytic continuation is needed for the coefficient calculation.
More explicitly, with row $r=N-2$ and column $n$,
\[
 (K_b)_{N-2,n-1}
 =h_{N-n}\{2n(2n-1)+(2N-1)(2N-2)-2\mathbf1_{n=1}\}
\]
when $n\le N$, and the entry is zero when $n>N$.

Write $Q_*(t)=\sum_{n\ge1}q_n^*t^{2n}$ for the formal solution
of $\mathcal L_{Q_*}=0$ normalized by $q_1^*=1$.  For $n\ge2$,
the coefficient of $t^{2n-3}$ gives
\begin{equation}\label{adj:q-recurrence}
 4(2n-1)^2q_n^*
 +\sum_{r=1}^{n-1}h_{n-r}
 \left\{2r(2r-1)+(2n-1)(2n-2)-2\mathbf1_{r=1}\right\}q_r^*=0.
\end{equation}
The nonzero coefficient of $q_n^*$ proves existence and uniqueness
recursively.  Equations \eqref{adj:hsign} and
\eqref{adj:positive-recurrence} show that
\begin{equation}\label{adj:qsign}
 q_n^*=(-1)^{n-1}c_n.
\end{equation}
In particular no coefficient vanishes.

The matrix $K_b$ is upper Hessenberg.  Its row corresponding to
$n=2,\ldots,b+1$ has the nonzero next-column entry
$d_n=4(2n-1)^2$ whenever that column is present.  Successively
using these entries to eliminate $q_2,\ldots,q_b$ shows
\begin{equation}\label{adj:hessenberg-det}
 \det K_b=(-1)^b q_{b+1}^*\prod_{n=2}^{b+1}d_n
 =4^b((2b+1)!!)^2c_{b+1}.
\end{equation}
Here the last constraint is the one that would determine
$q_{b+1}^*$ in the infinite formal solution; a polynomial has
$q_{b+1}=0$.  One explicit determinant argument replaces the first
column by itself plus $q_2^*,\ldots,q_b^*$ times the successive
columns.  Its first $b-1$ entries become zero and its last entry
is $-d_{b+1}q_{b+1}^*$.  Expanding along that column leaves a
lower triangular minor with diagonal $d_2,\ldots,d_b$, proving
\eqref{adj:hessenberg-det} with its stated sign.
This also provides a determinant-free proof of
injectivity: the first $b-1$ constraints force
$q_n=q_1q_n^*$ through $n=b$, and the final one forces
$q_1q_{b+1}^*=0$, hence $Q=0$.

Taking determinants in \eqref{adj:matrix-conjugacy}, and using
\[
 \prod_{n=1}^{b}\frac{(2n-1)!}{(2n)!}=\frac1{2^bb!},
\]
proves \eqref{adj:determinant}.  The transform
\eqref{adj:forward} is lower triangular with diagonal
$1/(2r+1)$, and the first $b$ block rows are the identity on
$C_0,\ldots,C_{b-1}$.  Its product of diagonal entries through
$r=2b-1$ is $1/(4b-1)!!$.  This proves
\eqref{adj:full-determinant} and the rank assertion.  Finally
$M_b^{-1}$ reconstructs every even-zeta coordinate from
$C_b,\ldots,C_{2b-1}$, while \eqref{adj:backward} obtains all
these blocks from the first $2b$ moments.  Equations
\eqref{adj:a} and \eqref{adj:forward} then reconstruct every
later moment by finite rational arithmetic.
\end{proof}

\paragraph{A clarification about the range of the theorem.}
The determinant formula proves the full infinite family $q-p=2$.
The case $p=q$ was already proved in the incoming report.  General
unequal gaps $q-p\ge4$ remain open here; the argument is not a
proof of the complete source conjecture.  Its positivity mechanism
uses a single free initial coefficient, which is specific to the
second-order operator in \eqref{adj:operator}.
Isolated larger-gap pairs, including $(p,q)=(2,6)$, were already
treated in the source; the remaining question is the uniform assertion.

\subsection{Ordinary convergent harmonic and polygamma identities}

The inverse matrix produces explicit series identities without an
unevaluated finite-part symbol.  For $0\le k<b$, define
\begin{align}
 w_{b,k,r}
 &=\sum_{t=0}^{b-1}(M_b^{-1})_{kt}
 \binom{2(b+t)+1}{2r}2^{-2(b+t-r)},
 \label{adj:w}\\
 P_{b,k}(x)&=\sum_{r=0}^{2b-1}w_{b,k,r}x^{2r},
 \label{adj:P}
\end{align}
where a summand in \eqref{adj:w} is zero if $r>b+t$.  Put
\begin{equation}\label{adj:alpha}
 a_0=\frac1{(2b-3)(2b-1)},\qquad
 a_1=-\frac{2b-2}{24(2b-1)}-\frac{2b}{24(2b-3)}.
\end{equation}

\begin{corollary}[All even-zeta coordinates from adjacent tails]
\label{adj:ordinary}
For every $b\ge2$ and $0\le k<b$,
\begin{equation}\label{adj:ordinary-family}
 \begin{split}
 \sum_{n=0}^{\infty}\bigl\{&P_{b,k}(x_n)A_{2b-2}(n)A_{2b}(n)
       -a_0w_{b,k,2b-1}x_n^2\\
 &\hspace{12mm}-a_0w_{b,k,2b-2}-a_1w_{b,k,2b-1}\bigr\}
 =\zeta(2k).
 \end{split}
\end{equation}
The series is absolutely convergent.  Replacing the tail product by
\begin{equation}\label{adj:polygamma}
 A_{2b-2}(n)A_{2b}(n)
 =\frac{\psi^{(2b-3)}(n+1)\psi^{(2b-1)}(n+1)}
 {(2b-3)!(2b-1)!}
\end{equation}
gives the corresponding ordinary polygamma-product identity.
\end{corollary}

\begin{proof}
By \eqref{adj:block}, inversion, and \eqref{adj:backward},
$\sum_rw_{b,k,r}D_r=\zeta(2k)$.  The centered Hurwitz
expansion gives
\[
 A_{2b-2}(n)A_{2b}(n)
 =a_0x_n^{4-4b}+a_1x_n^{2-4b}+O(x_n^{-4b}).
\]
The largest possible degree in \eqref{adj:P} is $4b-2$.
Thus the two displayed subtractions remove every nondecaying term,
and the summand is $O(x_n^{-2})$.  Centered subtraction computes
the same continued moment because the removed monomials have
$\zeta(0,1/2)=\zeta(-2,1/2)=0$.  This proves the identity and
absolute convergence.  Formula \eqref{adj:polygamma} is the
standard Hurwitz-zeta expression for odd-order polygammas.
\end{proof}

\subsection{Two exact rational evaluations}

For $(p,q)=(4,6)$, the late block matrix and its inverse are
\begin{equation}\label{adj:46matrix}
 M_3=\begin{pmatrix}
 0&3/4&1/4\\5/12&2/3&-1/12\\1&-3/8&1/12
 \end{pmatrix},\qquad
 M_3^{-1}=\frac1{339}
 \begin{pmatrix}-28&180&264\\136&288&-120\\948&-864&360\end{pmatrix}.
\end{equation}
In particular $\det M_3=-113/384$.  Define
\begin{equation}\label{adj:P46}
 \begin{split}
 P_{4,6}(x)={}&67+15516x^2+324320x^4+1434496x^6\\
             &+1601280x^8+371712x^{10}.
 \end{split}
\end{equation}
Then the first row of \eqref{adj:46matrix} after the inverse
Bernoulli transform gives the ordinary identity
\begin{equation}\label{adj:46rational}
 \sum_{n=0}^{\infty}\left\{
 P_{4,6}(x_n)A_4(n)A_6(n)
 -\frac{123904}{5}x_n^2-\frac{316928}{5}
 \right\}=-21696.
\end{equation}
The right side is $43392\zeta(0)$.  Thus every double-zeta
coordinate and every positive even-zeta coordinate cancels
exactly.

For $(p,q)=(6,8)$, $\det M_4=2509/5120$.  Define
\begin{equation}\label{adj:P68}
 \begin{split}
 P_{6,8}(x)={}&3705+880180x^2+32541904x^4+389667904x^6\\
 &+1652632320x^8+2454887424x^{10}\\
 &+1189933056x^{12}+138854400x^{14}.
 \end{split}
\end{equation}
Its rational-coordinate identity is
\begin{equation}\label{adj:68rational}
 \sum_{n=0}^{\infty}\left\{
 P_{6,8}(x_n)A_6(n)A_8(n)
 -\frac{27770880}{7}x_n^2-\frac{98910208}{5}
 \right\}=-7707648.
\end{equation}
The product in \eqref{adj:46rational} is
$\psi^{(3)}(n+1)\psi^{(5)}(n+1)/720$; the product in
\eqref{adj:68rational} is
$\psi^{(5)}(n+1)\psi^{(7)}(n+1)/604800$.
The coefficient files also contain the corresponding exact
$\zeta(2),\zeta(4)$, and $\zeta(6)$ rows where applicable.

\subsection{A hypergeometric expression for the determinant generator}

The formal solution used in the proof has a separate exact expression.
It is not needed for nonvanishing, but supplies a second coefficient
algorithm.  Dividing \eqref{adj:operator} by $2h(t)$ gives
\begin{equation}\label{adj:Qode}
 Q_*''(t)-\operatorname{csch}(t)Q_*'(t)
 +\frac{Q_*(t)}{4\sinh^2(t/2)}=1.
\end{equation}
Set $z=\tanh(t/2)$ and $Q_*(t)=4z^2F(z^2)$.  If
$F(u)=\sum_{n\ge0}f_nu^n$, then
\begin{equation}\label{adj:frecurrence}
 f_0=1,\qquad (n+\tfrac12)^2f_n-n^2f_{n-1}=\tfrac14
 \quad(n\ge1),
\end{equation}
and hence
\begin{equation}\label{adj:fclosed}
 f_n=\left(\frac{n!}{(3/2)_n}\right)^2
 \sum_{k=0}^{n}\left(\frac{(1/2)_k}{k!}\right)^2.
\end{equation}
This follows by dividing \eqref{adj:frecurrence} by its homogeneous
product and telescoping.  It is also a direct formal derivation of
\eqref{adj:Qode}; the series converges near the origin, since the
coefficients in \eqref{adj:fclosed} have at most polynomial growth.
Consequently
\begin{equation}\label{adj:detgenerator}
 \sum_{n=1}^{\infty}(-1)^{n-1}c_nt^{2n}
 =4\tanh^2(t/2)
 \sum_{n=0}^{\infty}f_n\tanh^{2n}(t/2)
\end{equation}
as an analytic identity near zero.  The finite binomial-square sum
in \eqref{adj:fclosed} connects the universal determinant
coefficients to the classical complete-elliptic-integral
hypergeometric series.  The present theorem uses this connection as
an alternative exact representation, without claiming a new
historical priority for that classical series.

\subsection{Checks, source audit, and the remaining conjecture}

The accompanying \texttt{verify\_adjacent\_tail\_prefix.py} compares
the original Bernoulli-binomial matrix with an independent
differential-operator construction, checks the determinant formula
through $b=16$, and independently reconstructs the two displayed
examples from the source $U/R$ block formulas, retaining formal
double-zeta coordinates until their exact cancellation.  Its
numerical option evaluates rational and nonrational rows by direct
tail sums followed by an independently generated centered-Hurwitz
tail.  Those residuals are diagnostics, not interval certificates
or substitutes for the proofs above.

No mathematical error in the inherited tail theorem is needed for
this advance.  Two convention errors are explicitly excluded:
omitting $Q''(0)h(t)$ changes the operator, and omitting the column
reversal changes the determinant sign.  The source optimal-prefix
conjecture should be updated to record the newly proved adjacent
family, while retaining open status for general gaps at least four.
